% !TEX TS-program = lualatex
% !TEX encoding = UTF-8 Unicode
\documentclass[12pt,a4paper]{article}

\IfFileExists{src/libs/body.tex}{% --- body.tex ---
\usepackage{pdfpages}
\usepackage{fontspec}
% Prefer Times New Roman when available; fall back in container environments.
\IfFontExistsTF{Times New Roman}{
  \setmainfont{Times New Roman}
}{
  \setmainfont{Liberation Serif}
}
\usepackage[vietnamese]{babel}

% Page layout and spacing
\usepackage[left=3cm,right=2cm,top=2cm,bottom=2cm]{geometry}
\usepackage{titlesec}
\usepackage{setspace}

% Silence noisy warnings
\usepackage{silence}
\WarningFilter{transparent}{Loading aborted}
\WarningFilter{hyperref}{The anchor of a bookmark and its parent's}

% Core packages
\usepackage{float}
\usepackage{svg}
\usepackage{graphicx}
\graphicspath{{assets/}{../assets/}{../../assets/}}
\svgpath{{assets/}{../assets/}{../../assets/}}
\usepackage{enumitem}
\usepackage{array}
\usepackage{tabularx}
\usepackage{multirow}
\usepackage[table]{xcolor}
\usepackage{ulem}
\usepackage{xurl}

% ====== THÊM 2 GÓI NÀY ĐỂ FIX LỖI ADJUSTBOX VÀ FOREST ======
\usepackage{adjustbox}
\usepackage[edges]{forest}
% ============================================================

% TikZ
\usepackage{tikz}

% ====== BỔ SUNG FIT, BACKGROUNDS, TREES VÀO DANH SÁCH ======
\usetikzlibrary{calc, arrows.meta, shapes.geometric, positioning, fit, backgrounds, trees}

% Table tuning
\renewcommand{\tabularxcolumn}[1]{m{#1}}
\hbadness=10000
\hfuzz=10000pt

% Shared commands
\newcommand{\subsecheading}[1]{%
    \par\vspace{2pt}%
    \hspace{1.5em}\textbf{#1}%
    \par\vspace{2pt}%
}
\newcommand{\introtext}[1]{%
    \par\hspace{3.0em}#1\par%
}
\newcommand{\StandaloneDocumentSetup}[2]{%
  \pagenumbering{arabic}%
  \ApplyStandardFooter{#1}{#2}%
}
\newcommand{\StandaloneSectionSetup}[2]{%
  \ApplyStandardFooter{#1}{#2}%
}

% Math
\usepackage{amsmath}

% Hyperref
\usepackage{hyperref}
\hypersetup{
    colorlinks=true,
    linkcolor=black,
    filecolor=magenta,
    urlcolor=blue,
}}{%
  \IfFileExists{../libs/body.tex}{% --- body.tex ---
\usepackage{pdfpages}
\usepackage{fontspec}
% Prefer Times New Roman when available; fall back in container environments.
\IfFontExistsTF{Times New Roman}{
  \setmainfont{Times New Roman}
}{
  \setmainfont{Liberation Serif}
}
\usepackage[vietnamese]{babel}

% Page layout and spacing
\usepackage[left=3cm,right=2cm,top=2cm,bottom=2cm]{geometry}
\usepackage{titlesec}
\usepackage{setspace}

% Silence noisy warnings
\usepackage{silence}
\WarningFilter{transparent}{Loading aborted}
\WarningFilter{hyperref}{The anchor of a bookmark and its parent's}

% Core packages
\usepackage{float}
\usepackage{svg}
\usepackage{graphicx}
\graphicspath{{assets/}{../assets/}{../../assets/}}
\svgpath{{assets/}{../assets/}{../../assets/}}
\usepackage{enumitem}
\usepackage{array}
\usepackage{tabularx}
\usepackage{multirow}
\usepackage[table]{xcolor}
\usepackage{ulem}
\usepackage{xurl}

% ====== THÊM 2 GÓI NÀY ĐỂ FIX LỖI ADJUSTBOX VÀ FOREST ======
\usepackage{adjustbox}
\usepackage[edges]{forest}
% ============================================================

% TikZ
\usepackage{tikz}

% ====== BỔ SUNG FIT, BACKGROUNDS, TREES VÀO DANH SÁCH ======
\usetikzlibrary{calc, arrows.meta, shapes.geometric, positioning, fit, backgrounds, trees}

% Table tuning
\renewcommand{\tabularxcolumn}[1]{m{#1}}
\hbadness=10000
\hfuzz=10000pt

% Shared commands
\newcommand{\subsecheading}[1]{%
    \par\vspace{2pt}%
    \hspace{1.5em}\textbf{#1}%
    \par\vspace{2pt}%
}
\newcommand{\introtext}[1]{%
    \par\hspace{3.0em}#1\par%
}
\newcommand{\StandaloneDocumentSetup}[2]{%
  \pagenumbering{arabic}%
  \ApplyStandardFooter{#1}{#2}%
}
\newcommand{\StandaloneSectionSetup}[2]{%
  \ApplyStandardFooter{#1}{#2}%
}

% Math
\usepackage{amsmath}

% Hyperref
\usepackage{hyperref}
\hypersetup{
    colorlinks=true,
    linkcolor=black,
    filecolor=magenta,
    urlcolor=blue,
}}{% --- body.tex ---
\usepackage{pdfpages}
\usepackage{fontspec}
% Prefer Times New Roman when available; fall back in container environments.
\IfFontExistsTF{Times New Roman}{
  \setmainfont{Times New Roman}
}{
  \setmainfont{Liberation Serif}
}
\usepackage[vietnamese]{babel}

% Page layout and spacing
\usepackage[left=3cm,right=2cm,top=2cm,bottom=2cm]{geometry}
\usepackage{titlesec}
\usepackage{setspace}

% Silence noisy warnings
\usepackage{silence}
\WarningFilter{transparent}{Loading aborted}
\WarningFilter{hyperref}{The anchor of a bookmark and its parent's}

% Core packages
\usepackage{float}
\usepackage{svg}
\usepackage{graphicx}
\graphicspath{{assets/}{../assets/}{../../assets/}}
\svgpath{{assets/}{../assets/}{../../assets/}}
\usepackage{enumitem}
\usepackage{array}
\usepackage{tabularx}
\usepackage{multirow}
\usepackage[table]{xcolor}
\usepackage{ulem}
\usepackage{xurl}

% ====== THÊM 2 GÓI NÀY ĐỂ FIX LỖI ADJUSTBOX VÀ FOREST ======
\usepackage{adjustbox}
\usepackage[edges]{forest}
% ============================================================

% TikZ
\usepackage{tikz}

% ====== BỔ SUNG FIT, BACKGROUNDS, TREES VÀO DANH SÁCH ======
\usetikzlibrary{calc, arrows.meta, shapes.geometric, positioning, fit, backgrounds, trees}

% Table tuning
\renewcommand{\tabularxcolumn}[1]{m{#1}}
\hbadness=10000
\hfuzz=10000pt

% Shared commands
\newcommand{\subsecheading}[1]{%
    \par\vspace{2pt}%
    \hspace{1.5em}\textbf{#1}%
    \par\vspace{2pt}%
}
\newcommand{\introtext}[1]{%
    \par\hspace{3.0em}#1\par%
}
\newcommand{\StandaloneDocumentSetup}[2]{%
  \pagenumbering{arabic}%
  \ApplyStandardFooter{#1}{#2}%
}
\newcommand{\StandaloneSectionSetup}[2]{%
  \ApplyStandardFooter{#1}{#2}%
}

% Math
\usepackage{amsmath}

% Hyperref
\usepackage{hyperref}
\hypersetup{
    colorlinks=true,
    linkcolor=black,
    filecolor=magenta,
    urlcolor=blue,
}}%
}
\IfFileExists{src/libs/header.tex}{% --- header.tex ---
\usepackage{fancyhdr}
\setlength{\headheight}{15.5pt} % Thêm dòng này để fix cảnh báo chiều cao
%\setlength{\headheight}{14.5pt}
\addtolength{\topmargin}{-2.5pt}

\pagestyle{fancy}
\fancyhf{}
\renewcommand{\headrulewidth}{0pt}
\renewcommand{\footrulewidth}{0pt}

\newcommand{\ApplyStandardHeader}{%
  \fancyhead[L]{%
    \begin{minipage}[t]{0.795\linewidth}%
      UIT \textbar{} SS004.F21 \textbar{} Nhóm 03 \textbar{} Đồ án cuối kỳ%
    \end{minipage}%
    \vrule%
    \begin{minipage}[t]{0.185\linewidth}%
      \centering cTetris%
    \end{minipage}%
  }%
}

\ApplyStandardHeader
}{%
  \IfFileExists{../libs/header.tex}{% --- header.tex ---
\usepackage{fancyhdr}
\setlength{\headheight}{15.5pt} % Thêm dòng này để fix cảnh báo chiều cao
%\setlength{\headheight}{14.5pt}
\addtolength{\topmargin}{-2.5pt}

\pagestyle{fancy}
\fancyhf{}
\renewcommand{\headrulewidth}{0pt}
\renewcommand{\footrulewidth}{0pt}

\newcommand{\ApplyStandardHeader}{%
  \fancyhead[L]{%
    \begin{minipage}[t]{0.795\linewidth}%
      UIT \textbar{} SS004.F21 \textbar{} Nhóm 03 \textbar{} Đồ án cuối kỳ%
    \end{minipage}%
    \vrule%
    \begin{minipage}[t]{0.185\linewidth}%
      \centering cTetris%
    \end{minipage}%
  }%
}

\ApplyStandardHeader
}{% --- header.tex ---
\usepackage{fancyhdr}
\setlength{\headheight}{15.5pt} % Thêm dòng này để fix cảnh báo chiều cao
%\setlength{\headheight}{14.5pt}
\addtolength{\topmargin}{-2.5pt}

\pagestyle{fancy}
\fancyhf{}
\renewcommand{\headrulewidth}{0pt}
\renewcommand{\footrulewidth}{0pt}

\newcommand{\ApplyStandardHeader}{%
  \fancyhead[L]{%
    \begin{minipage}[t]{0.795\linewidth}%
      UIT \textbar{} SS004.F21 \textbar{} Nhóm 03 \textbar{} Đồ án cuối kỳ%
    \end{minipage}%
    \vrule%
    \begin{minipage}[t]{0.185\linewidth}%
      \centering cTetris%
    \end{minipage}%
  }%
}

\ApplyStandardHeader
}%
}
\IfFileExists{src/libs/footer.tex}{% --- footer.tex ---
\usepackage{lastpage}

\newcommand{\ApplyStandardFooter}[2]{%
  \fancyfoot[L]{%
    \begin{minipage}[t]{0.795\linewidth}%
      Document ID: #1%
    \end{minipage}%
    \vrule%
    \begin{minipage}[t]{0.185\linewidth}%
      \centering\thepage/\pageref{#2}%
    \end{minipage}%
  }%
  \fancypagestyle{plain}{%
    \fancyhf{}%
    \renewcommand{\headrulewidth}{0pt}%
    \renewcommand{\footrulewidth}{0pt}%
    \ApplyStandardHeader%
    \fancyfoot[L]{%
      \begin{minipage}[t]{0.795\linewidth}%
        Document ID: #1%
      \end{minipage}%
      \vrule%
      \begin{minipage}[t]{0.185\linewidth}%
        \centering\thepage/\pageref{#2}%
      \end{minipage}%
    }%
  }%
  \pagestyle{fancy}%
}
}{%
  \IfFileExists{../libs/footer.tex}{% --- footer.tex ---
\usepackage{lastpage}

\newcommand{\ApplyStandardFooter}[2]{%
  \fancyfoot[L]{%
    \begin{minipage}[t]{0.795\linewidth}%
      Document ID: #1%
    \end{minipage}%
    \vrule%
    \begin{minipage}[t]{0.185\linewidth}%
      \centering\thepage/\pageref{#2}%
    \end{minipage}%
  }%
  \fancypagestyle{plain}{%
    \fancyhf{}%
    \renewcommand{\headrulewidth}{0pt}%
    \renewcommand{\footrulewidth}{0pt}%
    \ApplyStandardHeader%
    \fancyfoot[L]{%
      \begin{minipage}[t]{0.795\linewidth}%
        Document ID: #1%
      \end{minipage}%
      \vrule%
      \begin{minipage}[t]{0.185\linewidth}%
        \centering\thepage/\pageref{#2}%
      \end{minipage}%
    }%
  }%
  \pagestyle{fancy}%
}
}{% --- footer.tex ---
\usepackage{lastpage}

\newcommand{\ApplyStandardFooter}[2]{%
  \fancyfoot[L]{%
    \begin{minipage}[t]{0.795\linewidth}%
      Document ID: #1%
    \end{minipage}%
    \vrule%
    \begin{minipage}[t]{0.185\linewidth}%
      \centering\thepage/\pageref{#2}%
    \end{minipage}%
  }%
  \fancypagestyle{plain}{%
    \fancyhf{}%
    \renewcommand{\headrulewidth}{0pt}%
    \renewcommand{\footrulewidth}{0pt}%
    \ApplyStandardHeader%
    \fancyfoot[L]{%
      \begin{minipage}[t]{0.795\linewidth}%
        Document ID: #1%
      \end{minipage}%
      \vrule%
      \begin{minipage}[t]{0.185\linewidth}%
        \centering\thepage/\pageref{#2}%
      \end{minipage}%
    }%
  }%
  \pagestyle{fancy}%
}
}%
}
\usepackage{docmute}
\hypersetup{pdftitle={Báo cáo Kỹ năng nghề nghiệp - Danh sách kỹ năng}}

\begin{document}
\providecommand{\StandaloneSectionSetup}[2]{\ApplyStandardFooter{#1}{#2}}
\StandaloneSectionSetup{17-skillsetList.tex}{LastPageSkillsetList}
\onehalfspacing

\section{Danh sách kỹ năng nghề nghiệp}

Phần này phục vụ hai mục tiêu đồng thời: (1) \textbf{tuyển dụng} --- xác định rõ kỹ năng
cần thiết theo từng vai trò để lựa chọn ứng viên phù hợp; (2) \textbf{đào tạo tự học}
--- cung cấp lộ trình học tập có thứ tự ưu tiên để thành viên mới có thể tự onboarding
vào hệ sinh thái cTetris mà không cần hướng dẫn liên tục.

% ══════════════════════════════════════════════════════════════════
\subsection{Ma trận kỹ năng theo vai trò}
% ══════════════════════════════════════════════════════════════════

Dự án có năm vai trò. Bảng dưới định nghĩa yêu cầu tối thiểu (\levelBeg),
khuyến nghị (\levelMid) và lý tưởng (\levelAdv) để đóng góp hiệu quả ngay từ tuần đầu.
Ma trận được chia làm hai phần: \textbf{kỹ năng kỹ thuật} (công cụ, ngôn ngữ, framework)
và \textbf{kỹ năng mềm} (giao tiếp, quản lý, làm việc nhóm). Cả hai đều quan trọng như nhau
trong môi trường dự án phân tán đa nền tảng với năm thành viên thuộc các module khác nhau.

\vspace{4pt}
\noindent\textbf{a) Ma trận kỹ năng kỹ thuật}

\vspace{2pt}
{\small
\setlength{\tabcolsep}{4pt}
\begin{tabularx}{\linewidth}{|X|c|c|c|c|c|}
\hline
\rowcolor{gray!25}
\textbf{Kỹ năng} &
\rotatebox{60}{\shortstack[l]{\textbf{Quản lý}\\\textbf{dự án}}} &
\rotatebox{60}{\shortstack[l]{\textbf{Lập trình viên}\\\textbf{C++/SDL3}}} &
\rotatebox{60}{\shortstack[l]{\textbf{Lập trình viên}\\\textbf{C++/SQLite}}} &
\rotatebox{60}{\shortstack[l]{\textbf{Lập trình viên}\\\textbf{C++/Git}}} &
\rotatebox{60}{\shortstack[l]{\textbf{QA /}\\\textbf{Kiểm thử}}} \\
\hline
C++17 (OOP, STL, con trỏ)   & \levelBeg & \levelAdv & \levelAdv & \levelMid & \levelBeg \\
\hline
CMake và hệ thống build       & \levelMid & \levelMid & \levelMid & \levelMid & \levelBeg \\
\hline
SDL3 (renderer, event, audio) & \levelBeg & \levelAdv & \levelAdv & \levelMid & \levelBeg \\
\hline
SQLite và SQL cơ bản          & \levelBeg & \levelMid & \levelAdv & \levelMid & \levelBeg \\
\hline
Emscripten / WebAssembly      & \levelBeg & \levelMid & \levelMid & \levelBeg & \levelBeg \\
\hline
Git và GitHub workflow         & \levelAdv & \levelAdv & \levelAdv & \levelAdv & \levelMid \\
\hline
JSON và REST API               & \levelBeg & \levelBeg & \levelMid & \levelAdv & \levelBeg \\
\hline
LaTeX / TikZ                  & \levelAdv & \levelMid & \levelMid & \levelMid & \levelAdv \\
\hline
Kiểm thử phần mềm             & \levelMid & \levelMid & \levelMid & \levelMid & \levelAdv \\
\hline
Thiết kế narrative / storyline & \levelBeg & \levelBeg & \levelBeg & \levelAdv & \levelBeg \\
\hline
\end{tabularx}
}

\vspace{8pt}
\noindent\textbf{b) Ma trận kỹ năng mềm}

\vspace{2pt}
{\small
\setlength{\tabcolsep}{4pt}
\begin{tabularx}{\linewidth}{|X|c|c|c|c|c|}
\hline
\rowcolor{gray!25}
\textbf{Kỹ năng} &
\rotatebox{60}{\shortstack[l]{\textbf{Quản lý}\\\textbf{dự án}}} &
\rotatebox{60}{\shortstack[l]{\textbf{Lập trình viên}\\\textbf{C++/SDL3}}} &
\rotatebox{60}{\shortstack[l]{\textbf{Lập trình viên}\\\textbf{C++/SQLite}}} &
\rotatebox{60}{\shortstack[l]{\textbf{Lập trình viên}\\\textbf{C++/Git}}} &
\rotatebox{60}{\shortstack[l]{\textbf{QA /}\\\textbf{Kiểm thử}}} \\
\hline
Giao tiếp bằng văn bản (Slack, PR, email)
                              & \levelAdv & \levelMid & \levelMid & \levelAdv & \levelAdv \\
\hline
Giao tiếp trực tiếp (họp, thuyết trình)
                              & \levelAdv & \levelMid & \levelMid & \levelMid & \levelMid \\
\hline
Lập kế hoạch và quản lý thời gian
                              & \levelAdv & \levelMid & \levelMid & \levelMid & \levelMid \\
\hline
Quản lý xung đột và đàm phán
                              & \levelAdv & \levelBeg & \levelBeg & \levelMid & \levelMid \\
\hline
Lãnh đạo và uỷ thác công việc
                              & \levelAdv & \levelBeg & \levelBeg & \levelBeg & \levelMid \\
\hline
Làm việc nhóm và phối hợp đa module
                              & \levelAdv & \levelAdv & \levelAdv & \levelAdv & \levelAdv \\
\hline
Tư duy phản biện và đặt câu hỏi đúng
                              & \levelMid & \levelMid & \levelMid & \levelMid & \levelAdv \\
\hline
Phản hồi xây dựng (code review, feedback)
                              & \levelAdv & \levelAdv & \levelAdv & \levelAdv & \levelAdv \\
\hline
Tài liệu hoá và chia sẻ tri thức
                              & \levelAdv & \levelMid & \levelMid & \levelAdv & \levelAdv \\
\hline
Học tập chủ động và thích ứng
                              & \levelMid & \levelMid & \levelMid & \levelMid & \levelMid \\
\hline
Agile / Kanban (Trello, Slack)
                              & \levelAdv & \levelMid & \levelMid & \levelMid & \levelMid \\
\hline
\end{tabularx}
}

% ══════════════════════════════════════════════════════════════════
\subsection{Tiêu chí tuyển dụng theo vai trò}
% ══════════════════════════════════════════════════════════════════

\subsubsection{Chuyên viên quản lý dự án (Project Manager)}

\noindent\textbf{Yêu cầu bắt buộc:}
\begin{itemize}[nosep]
    \item Thành thạo Git workflow: nhánh, merge, conflict resolution, Pull Request.
    \item Sử dụng được Trello (Kanban) và Slack để theo dõi và thông báo tiến độ.
    \item Đọc hiểu được tài liệu LaTeX và phát hiện sai sót cấu trúc file \texttt{.tex}.
    \item Có khả năng tổng hợp, chuẩn hóa báo cáo và mã nguồn từ nhiều thành viên.
\end{itemize}

\noindent\textbf{Ưu tiên thêm:}
\begin{itemize}[nosep]
    \item Hiểu biết về C++17 ở mức đọc và nhận xét code (không cần viết thành thạo).
    \item Kinh nghiệm quản lý dự án nhóm 3--6 người theo phương pháp Agile.
    \item Biết dùng GitHub Actions để theo dõi kết quả CI/CD.
\end{itemize}

\noindent\textbf{Kỹ năng mềm cốt lõi:}
\begin{itemize}[nosep]
    \item \textit{Giao tiếp đa kênh:} truyền đạt rõ ràng qua văn bản (Slack, email tổng hợp tuần) và trực tiếp (họp standup, họp review); biết chọn kênh phù hợp với mức độ khẩn cấp.
    \item \textit{Lãnh đạo phục vụ (servant leadership):} gỡ rào cản kỹ thuật/hành chính cho thành viên thay vì giao việc rồi bỏ mặc; chủ động hỏi ``Bạn cần hỗ trợ gì?'' trong mỗi standup.
    \item \textit{Quản lý xung đột:} can thiệp sớm khi có tranh luận về thiết kế hoặc quyền sở hữu code; áp dụng nguyên tắc ``tách vấn đề ra khỏi cá nhân''.
    \item \textit{Ra quyết định có dữ liệu:} dựa trên metric Trello (velocity, blocker time) thay vì cảm tính khi tái phân bổ task.
    \item \textit{Tư duy hệ thống:} nhìn ra phụ thuộc cross-module (vd: thay đổi \texttt{SettingsConfig} ảnh hưởng cả 3 module) và phối hợp release theo thứ tự đúng.
\end{itemize}

\subsubsection{Lập trình viên module gameCore (C++/SDL3) (Game Physics Developer)}

\noindent\textbf{Yêu cầu bắt buộc:}
\begin{itemize}[nosep]
    \item C++17 thành thạo: con trỏ, tham chiếu, struct/class, vòng lặp game loop.
    \item Hiểu vật lý bảng cờ Tetris: sinh khối, phát hiện va chạm, xóa dòng, tính điểm.
    \item SDL3: \texttt{SDL\_RenderFillRect}, \texttt{SDL\_PollEvent}, \texttt{SDL\_AudioStream}.
    \item Thuật toán: ma trận 2D, xoay toạ độ, kiểm tra biên (boundary check).
\end{itemize}

\noindent\textbf{Ưu tiên thêm:}
\begin{itemize}[nosep]
    \item Kinh nghiệm viết game loop 60 FPS không bị drop frame.
    \item Biết dùng Emscripten để build WASM từ C++ native.
    \item Hiểu SQLite ở mức viết \texttt{INSERT}, \texttt{SELECT} qua C API.
\end{itemize}

\noindent\textbf{Kỹ năng mềm cốt lõi:}
\begin{itemize}[nosep]
    \item \textit{Tài liệu hoá thuật toán:} viết comment và doc explain công thức xoay khối,
          chấm điểm để QA và đồng nghiệp hiểu mà không cần hỏi lại.
    \item \textit{Tiếp nhận feedback:} sẵn sàng chấp nhận sửa code khi QA phát hiện bug
          biên (boundary case) không phán xét người báo lỗi.
    \item \textit{Giao tiếp kỹ thuật:} mô tả bug và đề xuất giải pháp trong PR description
          theo cấu trúc: \textit{Vấn đề $\to$ Nguyên nhân $\to$ Giải pháp $\to$ Cách test}.
    \item \textit{Ước lượng công việc:} tự đánh giá khối lượng task theo Man-Day chính xác
          ($\pm$20\%); báo sớm nếu phát sinh khó khăn ngoài dự đoán.
\end{itemize}

\subsubsection{Lập trình viên module gameConsole (C++/SQLite) (UI / Hub Screen Developer)}

\noindent\textbf{Yêu cầu bắt buộc:}
\begin{itemize}[nosep]
    \item C++17 và SDL3 thành thạo; xây dựng UI procedural (không dùng framework UI).
    \item SQLite: thiết kế schema đa bảng, \texttt{CREATE TABLE IF NOT EXISTS}, UPSERT, JOIN.
    \item Xây dựng component tương tác: scrollbar, lightbox popup, focus navigation.
    \item Hiểu cross-platform: phát hiện và xử lý sai biệt giữa macOS, Windows, Ubuntu.
\end{itemize}

\noindent\textbf{Ưu tiên thêm:}
\begin{itemize}[nosep]
    \item Kinh nghiệm với Emscripten IDBFS (IndexedDB filesystem cho WASM).
    \item Biết tích hợp REST API (Cloudflare Workers) từ C++ (libcurl hoặc emscripten\_fetch).
    \item Có khả năng thiết kế hệ thống DB: 3 nhóm bảng, cascade unlock, foreign key logic.
\end{itemize}

\noindent\textbf{Kỹ năng mềm cốt lõi:}
\begin{itemize}[nosep]
    \item \textit{Đồng cảm với người dùng cuối (UX empathy):} hiểu rằng nút bấm khó nhìn,
          tooltip thiếu, hoặc focus navigation rối là ``bug trải nghiệm'' chứ không chỉ
          ``đẹp/xấu''.
    \item \textit{Phối hợp đa bên:} thường xuyên trao đổi với cả gameCore (logic game) và
          gameStory (nội dung) vì gameConsole là module trung gian; chủ động đặt câu hỏi
          làm rõ trước khi viết code thay vì giả định.
    \item \textit{Kiểm thử chéo nền tảng:} biên bản kiểm thử rõ ràng khi report bug
          cross-platform (macOS vs Windows vs Ubuntu): nêu cấu hình, phiên bản, các bước
          tái hiện.
    \item \textit{Kiên nhẫn với bug khó tái hiện:} không đổ lỗi cho người báo bug khi
          bug chỉ xảy ra ở 1 OS; tìm hiểu nguyên nhân hệ thống thay vì đóng ticket sớm.
\end{itemize}

\subsubsection{Lập trình viên module gameStory (C++/Git) (Content Sync \& Dialogue Developer)}

\noindent\textbf{Yêu cầu bắt buộc:}
\begin{itemize}[nosep]
    \item C++ ở mức trung cấp; hiểu cách parse JSON với nlohmann/json.
    \item Hiểu giao thức HTTP: GET request, response parsing, xử lý lỗi mạng.
    \item SQLite: đọc/ghi bảng nội dung (dialogue nodes, choices, chapter metadata).
    \item Python cơ bản: chỉnh sửa script \texttt{parse.py} và \texttt{manifest.py}.
\end{itemize}

\noindent\textbf{Ưu tiên thêm:}
\begin{itemize}[nosep]
    \item Kỹ năng viết nội dung narrative: xây dựng hội thoại phân nhánh, có chiều sâu cảm xúc.
    \item Hiểu GitHub Actions: chỉnh sửa file YAML workflow \texttt{sync-chapters.yml}.
    \item Biết SHA-diff sync: so sánh hash để chỉ tải về nội dung thay đổi.
\end{itemize}

\noindent\textbf{Kỹ năng mềm cốt lõi:}
\begin{itemize}[nosep]
    \item \textit{Tư duy biên tập (editorial mindset):} đọc và phản biện nội dung dialogue
          do người khác viết; đề xuất chỉnh sửa mà không làm mất giọng văn gốc.
    \item \textit{Giao tiếp với non-tech contributor:} hướng dẫn người viết nội dung
          (writer) cách commit lên Gist mà không cần biết Git command-line phức tạp.
    \item \textit{Tài liệu hoá quy trình:} viết hướng dẫn ``Cách thêm chapter mới'' đủ rõ
          để member mới làm theo được không cần hỏi.
    \item \textit{Quản lý phụ thuộc bên ngoài:} chủ động báo trước khi schema bảng dialogue
          thay đổi, vì điều này ảnh hưởng cả gameConsole và gameCore.
    \item \textit{Cân bằng nghệ thuật và kỹ thuật:} biết khi nào nên giữ nguyên ý đồ writer,
          khi nào cần đề xuất chỉnh sửa vì hạn chế kỹ thuật (vd: độ dài câu, encoding).
\end{itemize}

\subsubsection{Chuyên viên kiểm thử chất lượng phần mềm (QA Engineer)}

\noindent\textbf{Yêu cầu bắt buộc:}
\begin{itemize}[nosep]
    \item Đọc hiểu code C++ để viết test case có preconditions và expected output chính xác.
    \item Soạn test case thủ công theo mẫu: mã, tiêu đề, điều kiện tiên quyết, bước, kết quả.
    \item Kiểm thử cross-platform: chạy và so sánh hành vi trên macOS, Windows, Ubuntu.
    \item Viết biên bản họp, tổng hợp bug report theo cấu trúc chuẩn.
\end{itemize}

\noindent\textbf{Ưu tiên thêm:}
\begin{itemize}[nosep]
    \item Biết dùng Google Test hoặc Catch2 để viết unit test tự động.
    \item Hiểu CI/CD: đọc log GitHub Actions, phân biệt lỗi build và lỗi test.
    \item Có kinh nghiệm kiểm thử game: kiểm tra frame rate, input latency, memory leak.
\end{itemize}

\noindent\textbf{Kỹ năng mềm cốt lõi:}
\begin{itemize}[nosep]
    \item \textit{Viết bug report chuyên nghiệp:} cấu trúc \textit{Tiêu đề $\to$ Severity
          $\to$ Bước tái hiện $\to$ Hành vi mong đợi vs thực tế $\to$ Môi trường}; đính kèm
          screenshot/log để dev không phải hỏi lại.
    \item \textit{Giao tiếp không đối đầu (non-confrontational):} báo bug bằng cách mô tả
          hành vi quan sát được, không công kích người viết code (vd: ``Tính năng X không
          hoạt động khi Y'' thay vì ``Code của bạn sai'').
    \item \textit{Tư duy phản biện và đặt câu hỏi:} không tin tuyệt đối vào spec; thử các
          edge case mà developer có thể chưa nghĩ tới (input rỗng, overflow, race condition).
    \item \textit{Ưu tiên hoá bug:} phân loại Critical / Major / Minor / Trivial dựa trên
          tác động người dùng, không dựa trên độ khó kỹ thuật để fix.
    \item \textit{Kiên trì và tỉ mỉ:} thử lại bug nhiều lần với điều kiện khác nhau trước
          khi kết luận; không bỏ qua bug ``ngẫu nhiên'' --- chúng thường là race condition.
    \item \textit{Viết biên bản họp khách quan:} ghi quyết định và action item, không ghi
          chuyện cá nhân hay diễn giải chủ quan.
\end{itemize}

% ══════════════════════════════════════════════════════════════════
\subsection{Lộ trình tự học và onboarding cho thành viên mới}
% ══════════════════════════════════════════════════════════════════

Lộ trình được thiết kế theo thứ tự ưu tiên: kỹ năng ở mức cao hơn phụ thuộc vào
kỹ năng bên dưới. Thành viên mới nên hoàn thành từng bước trước khi tiếp theo.

\subsubsection{Giai đoạn 1 --- Cài đặt môi trường (Ngày 1--2)}

\begin{enumerate}[nosep]
    \item \textbf{Cài Git và clone repo:}
          \texttt{git clone https://github.com/phuctanpham/ctetris}
    \item \textbf{Cài trình biên dịch C++17:}
          macOS: \texttt{xcode-select --install};
          Ubuntu: \texttt{sudo apt install g++ cmake};
          Windows: cài MinGW-w64 hoặc Visual Studio 2022.
    \item \textbf{Chạy build lần đầu:}
          \texttt{bash build.sh} (macOS/Linux) hoặc \texttt{pwsh build.ps1} (Windows).
          Script tự phát hiện thiếu công cụ và in hướng dẫn cài.
    \item \textbf{Cài VS Code và extension:} C/C++ (Microsoft), CMake Tools, GitLens.
    \item \textbf{Tham gia Slack} và đọc kênh \texttt{\#general}, \texttt{\#dev-log}.
\end{enumerate}

\textbf{Kiểm tra:} chạy được module standalone (\texttt{BUILD\_STANDALONE=ON}) của
ít nhất một trong ba module: \texttt{gameStory}, \texttt{gameConsole}, \texttt{gameCore}.

\subsubsection{Giai đoạn 2 --- Hiểu cấu trúc dự án (Ngày 3--5)}

\begin{enumerate}[nosep]
    \item \textbf{Đọc tài liệu SRS} (§02, §03, §04) để hiểu yêu cầu nghiệp vụ từng module.
    \item \textbf{Đọc tài liệu SDS} (§05, §06, §07) để hiểu kiến trúc kỹ thuật, luồng dữ liệu.
    \item \textbf{Đọc \texttt{app/appTree.md}} trong repo để hiểu cây phụ thuộc module.
    \item \textbf{Chạy integrated build} (cả 3 module) và chơi game một vài ván.
    \item \textbf{Đọc hợp đồng nhóm} (§01-teamContract) và quy tắc làm việc.
\end{enumerate}

\textbf{Kiểm tra:} có thể giải thích luồng dữ liệu: \texttt{gameStory → gameConsole → gameCore}
và vai trò của \texttt{SettingsConfig}.

\subsubsection{Giai đoạn 3 --- Làm quen codebase (Ngày 6--10)}

\begin{tabularx}{\linewidth}{|l|X|l|}
\hline
\rowcolor{gray!20}
\textbf{File cần đọc} & \textbf{Kiến thức thu được} & \textbf{Ưu tiên} \\
\hline
\texttt{src/gameCore/app.cpp} & Game loop, GameState, spawnBlock, physics & \textbf{Cao} \\
\hline
\texttt{src/gameConsole/app.cpp} & AppState flags, DBLayer, SortEngine & \textbf{Cao} \\
\hline
\texttt{src/gameStory/app.cpp} & dialRunLoop, syncFromGist, postSyncCheck & Trung bình \\
\hline
\texttt{src/shared/gameConsole\_layout.h} & Struct SettingsConfig (cross-module contract) & \textbf{Cao} \\
\hline
\texttt{src/gameConsole/gameConsole\_db.h} & DB API: dbOpen, dbInsertRecord, dbClose & Trung bình \\
\hline
\texttt{CMakeLists.txt} & Cách tổ chức build, cờ \texttt{BUILD\_WASM}, link SDL3 & Trung bình \\
\hline
\end{tabularx}

\vspace{4pt}
\textbf{Kiểm tra:} có thể sửa một bug nhỏ (ví dụ: thay đổi màu nút, sửa công thức điểm) và
tạo Pull Request đúng quy trình.

\subsubsection{Giai đoạn 4 --- Kỹ năng chuyên sâu theo vai trò (Tuần 2--4)}

\noindent\textbf{Tài nguyên học tập tự học:}

\begin{tabularx}{\linewidth}{|l|X|l|}
\hline
\rowcolor{gray!20}
\textbf{Kỹ năng} & \textbf{Tài nguyên khuyến nghị} & \textbf{Thời lượng} \\
\hline
C++17 nâng cao & \url{https://learncpp.com} (chương 11--21) & 1--2 tuần \\
\hline
SDL3 & \url{https://wiki.libsdl.org/SDL3} (Getting Started) & 3--5 ngày \\
\hline
CMake & \url{https://cmake.org/cmake/help/latest/guide/tutorial} & 2--3 ngày \\
\hline
SQLite C API & \url{https://www.sqlite.org/quickstart.html} & 2--3 ngày \\
\hline
Emscripten / WASM & \url{https://emscripten.org/docs/getting_started} & 3--5 ngày \\
\hline
Git nâng cao & \url{https://learngitbranching.js.org} & 1--2 ngày \\
\hline
Google Test & \url{https://google.github.io/googletest} & 2--3 ngày \\
\hline
LaTeX / TikZ & \url{https://overleaf.com/learn/latex/Tutorials} & 1--2 tuần \\
\hline
GitHub Actions & \url{https://docs.github.com/en/actions/quickstart} & 1--2 ngày \\
\hline
Cloudflare Workers & \url{https://developers.cloudflare.com/workers} & 2--3 ngày \\
\hline
\end{tabularx}

\vspace{6pt}
\noindent\textbf{Tài nguyên học tập kỹ năng mềm:} (song song với kỹ năng kỹ thuật)

\vspace{2pt}
\begin{tabularx}{\linewidth}{|l|X|l|}
\hline
\rowcolor{gray!20}
\textbf{Kỹ năng} & \textbf{Tài nguyên khuyến nghị} & \textbf{Thời lượng} \\
\hline
Viết PR và commit message hiệu quả
    & \url{https://cbea.ms/git-commit/} (How to Write a Git Commit Message)
    & 1 ngày \\
\hline
Code review xây dựng
    & \url{https://google.github.io/eng-practices/review/} (Google Engineering Practices)
    & 2--3 ngày \\
\hline
Viết bug report chuyên nghiệp
    & \url{https://www.atlassian.com/agile/project-management/bug-tracking}
    & 1 ngày \\
\hline
Giao tiếp không đồng bộ (async communication)
    & GitLab Handbook: \url{https://handbook.gitlab.com/handbook/communication/}
    & 2--3 ngày \\
\hline
Quản lý dự án Agile / Scrum
    & \url{https://www.scrum.org/resources/scrum-guide} (Scrum Guide chính thức)
    & 1--2 ngày \\
\hline
Ước lượng và lập kế hoạch
    & \textit{Software Estimation: Demystifying the Black Art} --- Steve McConnell (sách)
    & 1 tuần \\
\hline
Giao tiếp kỹ thuật và viết tài liệu
    & \url{https://developers.google.com/tech-writing} (Google Tech Writing One)
    & 3--5 ngày \\
\hline
Quản lý xung đột nhóm
    & \textit{Crucial Conversations} --- Patterson et al. (sách, chương 1--6)
    & 1--2 tuần \\
\hline
Lãnh đạo kỹ thuật
    & \textit{The Manager's Path} --- Camille Fournier (chương 1--3 cho mọi vai trò)
    & 1 tuần \\
\hline
\end{tabularx}

\subsubsection{Giai đoạn 5 --- Đóng góp độc lập (Từ tuần 3 trở đi)}

\begin{enumerate}[nosep]
    \item Nhận task từ Trello với độ phức tạp phù hợp (bắt đầu từ \textit{Fix} hoặc \textit{Small Task}).
    \item Áp dụng mô hình \textbf{Design (0.5 MD) → Implement (1 MD) → Self-test (0.5 MD)}.
    \item Tạo Pull Request; đợi reviewer trong vòng 24 giờ; sửa theo feedback.
    \item Sau 2--3 PR được merge: có thể nhận \textit{Big Task} (nhóm 3+ hàm).
\end{enumerate}

% ══════════════════════════════════════════════════════════════════
\subsection{Kỹ năng mềm và văn hóa làm việc nhóm}
% ══════════════════════════════════════════════════════════════════

Trong một dự án phân tán với năm thành viên thuộc các module khác nhau, làm việc
trên ba hệ điều hành (macOS, Windows, Ubuntu), kỹ năng mềm có tác động trực tiếp
tới chất lượng sản phẩm không kém kỹ năng kỹ thuật. Phần này phân chia kỹ năng
mềm thành bốn nhóm: \textbf{giao tiếp}, \textbf{quản lý}, \textbf{lãnh đạo và
làm việc nhóm}, và \textbf{tự quản lý bản thân}.

\subsubsection{Nhóm 1 --- Kỹ năng giao tiếp}

\begin{tabularx}{\linewidth}{|l|X|c|}
\hline
\rowcolor{gray!20}
\textbf{Kỹ năng} & \textbf{Biểu hiện cụ thể trong dự án} & \textbf{Yêu cầu} \\
\hline
Giao tiếp chủ động & Phản hồi Slack trong 4 giờ làm việc; cập nhật Trello hàng ngày;
                     không để task ở trạng thái ``In Progress'' quá 2 ngày không có note. & Bắt buộc \\
\hline
Viết PR description & Cấu trúc: \textit{Mô tả thay đổi $\to$ Lý do $\to$ Cách test $\to$
                      Module bị ảnh hưởng}; tránh PR description một dòng. & Bắt buộc \\
\hline
Viết commit message & Theo Conventional Commits: \texttt{feat:}, \texttt{fix:},
                      \texttt{docs:}, \texttt{refactor:}; tiêu đề $\le$ 50 ký tự. & Bắt buộc \\
\hline
Họp hiệu quả & Đến đúng giờ; có agenda trước; ghi action item kèm người chịu trách
               nhiệm và hạn chót sau mỗi cuộc họp. & Bắt buộc \\
\hline
Thuyết trình kỹ thuật & Trình bày demo hoặc review module $\le$ 10 phút; nêu được
                        \textit{vấn đề, giải pháp, đánh đổi (trade-off)}. & Khuyến nghị \\
\hline
Giao tiếp đa văn hóa & Tôn trọng sự khác biệt khi nhóm có thành viên từ nhiều vùng;
                       dùng tiếng Việt chuẩn trong tài liệu chính thức, English trong code. & Khuyến nghị \\
\hline
Lắng nghe chủ động & Không ngắt lời khi người khác trình bày; xác nhận hiểu đúng bằng
                     cách tóm tắt lại (``Ý bạn là\ldots'') trước khi phản biện. & Bắt buộc \\
\hline
\end{tabularx}

\subsubsection{Nhóm 2 --- Kỹ năng quản lý (cho mọi vai trò, không riêng PM)}

\begin{tabularx}{\linewidth}{|l|X|c|}
\hline
\rowcolor{gray!20}
\textbf{Kỹ năng} & \textbf{Biểu hiện cụ thể trong dự án} & \textbf{Yêu cầu} \\
\hline
Quản lý thời gian cá nhân & Ước lượng Man-Day cho task $\pm$20\%; ưu tiên task
                            chặn người khác (blocker) trước task riêng. & Bắt buộc \\
\hline
Lập kế hoạch task & Phân nhỏ Big Task thành sub-task $\le$ 1 MD trước khi bắt đầu code;
                    không nhận task chưa rõ scope. & Bắt buộc \\
\hline
Cam kết deadline & Thông báo sớm tối thiểu 24 giờ nếu không kịp hạn; đề xuất kế hoạch
                   bù đắp (giảm scope, xin thêm thời gian, hoặc nhờ hỗ trợ). & Bắt buộc \\
\hline
Quản lý rủi ro & Nhận diện rủi ro kỹ thuật sớm (vd: thư viện không có version WASM, OS
                 không hỗ trợ API); báo cáo PM thay vì giấu để gỡ một mình. & Bắt buộc \\
\hline
Ưu tiên hoá (prioritization) & Phân biệt được \textit{quan trọng vs khẩn cấp};
                                bug Critical chặn release phải gỡ trước bug Minor. & Bắt buộc \\
\hline
Báo cáo tiến độ & Cập nhật Trello card với note ngắn cuối mỗi ngày làm việc;
                  PM tổng hợp được không cần hỏi từng người. & Bắt buộc \\
\hline
Quản lý phụ thuộc cross-module & Khi sửa \texttt{SettingsConfig} hoặc DB schema, báo
                                  trước 24h cho các module bị ảnh hưởng. & Bắt buộc \\
\hline
\end{tabularx}

\subsubsection{Nhóm 3 --- Lãnh đạo và làm việc nhóm}

\begin{tabularx}{\linewidth}{|l|X|c|}
\hline
\rowcolor{gray!20}
\textbf{Kỹ năng} & \textbf{Biểu hiện cụ thể trong dự án} & \textbf{Yêu cầu} \\
\hline
Phản hồi xây dựng & Review code tập trung vào logic, quy ước và khả năng bảo trì;
                    không phán xét cá nhân; nêu \textit{cái gì} và \textit{tại sao}. & Bắt buộc \\
\hline
Tiếp nhận phản hồi & Không phòng thủ khi bị review; coi feedback là cải tiến sản phẩm
                     chứ không phải tấn công cá nhân. & Bắt buộc \\
\hline
Tôn trọng quyền sở hữu & Không tự ý sửa file của người khác khi chưa có PR được duyệt;
                         báo trước nếu cần đụng tới phần code họ phụ trách. & Bắt buộc \\
\hline
Mentoring & Thành viên có kinh nghiệm sẵn lòng dành 15--30 phút/tuần hỗ trợ thành
            viên mới mà không xem là việc phụ. & Khuyến nghị \\
\hline
Quản lý xung đột & Khi có bất đồng kỹ thuật, đem ra thảo luận công khai trên kênh
                   \texttt{\#dev-discussion}; tránh nói xấu sau lưng; nếu không thống nhất
                   được, PM ra quyết định cuối. & Bắt buộc \\
\hline
Chia sẻ kiến thức & Khi học được kỹ thuật mới (vd: cách debug WASM), post lên kênh
                    \texttt{\#learning} với 2--3 dòng tóm tắt. & Khuyến nghị \\
\hline
Phối hợp đa module & Chủ động ping module liên quan khi thay đổi của mình có thể ảnh
                     hưởng tới họ; không đợi tới khi build vỡ. & Bắt buộc \\
\hline
Tư duy ``win-win'' & Khi review hoặc thảo luận, tìm giải pháp giúp cả hai bên cùng có lợi,
                     không phải tranh ai đúng ai sai. & Khuyến nghị \\
\hline
\end{tabularx}

\subsubsection{Nhóm 4 --- Tự quản lý bản thân}

\begin{tabularx}{\linewidth}{|l|X|c|}
\hline
\rowcolor{gray!20}
\textbf{Kỹ năng} & \textbf{Biểu hiện cụ thể trong dự án} & \textbf{Yêu cầu} \\
\hline
Tư duy sản phẩm & Mọi tính năng đều được hỏi: ``Người chơi được gì từ điều này?''
                  trước khi viết code. & Khuyến nghị \\
\hline
Tự học liên tục & Chủ động tra cứu tài liệu, StackOverflow, hoặc dùng AI assistant
                  trước khi hỏi đồng nghiệp; chia sẻ bài học với nhóm. & Khuyến nghị \\
\hline
Khả năng thích ứng & Sẵn sàng đổi tool, đổi cách làm khi nhóm thống nhất phương án mới;
                     không cứng nhắc theo thói quen cá nhân. & Bắt buộc \\
\hline
Tỉ mỉ và cẩn thận & Đọc lại code và test trước khi gửi PR; không gửi code chưa biên
                    dịch hay chưa chạy thử. & Bắt buộc \\
\hline
Đạo đức nghề nghiệp & Không copy code có bản quyền không phù hợp; trích nguồn khi
                      dùng snippet từ StackOverflow hoặc AI. & Bắt buộc \\
\hline
Quản lý năng lượng & Biết khi nào cần nghỉ; tránh code ``zombie'' lúc mệt vì dễ tạo bug
                     khó tìm hơn cả bug ban đầu. & Khuyến nghị \\
\hline
Trách nhiệm với lỗi & Khi gây bug regression, chủ động nhận trách nhiệm và sửa thay vì
                      đổ lỗi cho thiếu test hay tài liệu. & Bắt buộc \\
\hline
\end{tabularx}

\label{LastPageSkillsetList}
\end{document}



